\chapter{Workshop 1: A job worth monitoring}\label{workshop-01}

\section{Start with the user's question}

A person starts a build and wants to know whether it is making progress, whether it failed, and whether it is safe to leave. The job monitor in these workshops gives those questions a concrete home. It has three jobs: compile the application, run the tests, and publish the package. These are simulated jobs driven by scheduled events. The example never invokes a build tool or publishes anything. That boundary keeps the lesson about terminal application design while leaving room for a real process adapter later.

\subsection{Facts before decoration}

A job has a stable identity, a title, a textual status, a percentage and a run identity. Failure can also carry an explanation. The run identity distinguishes two attempts at the same job; a percentage alone cannot do that. Selection belongs to the monitor, rather than being inferred from where a highlight happened to be painted.

The \texttt{Job} class checks its invariants when a value is created. A successful job must be complete. Progress must be a finite number between zero and one hundred. A run identity is a non-negative integer. An immutable job can be handed to a renderer or retained in a test without a later update silently changing the value under examination.

\VerbatimInput[formatcom=\sourcecodefont,fontsize=\fontsize{12}{13}\selectfont,breaklines,breakanywhere]{companion-snippets/job.ts.txt}

\subsection{What immutability costs}

Each transition constructs a new job or monitor snapshot. That adds allocation and some copying. For three small jobs, the cost is easy to understand and measure. For a monitor with hundreds of thousands of entries, the collection strategy would need another decision. Immutability establishes an ownership boundary; it does not exempt a program from measuring its workload.

\section{Keep meaning out of the presentation}

A failed job remains failed when the terminal loses color. A running job remains running when a narrow viewport requires a different layout. The domain model therefore stores statuses such as \texttt{running} and \texttt{failed}, not color names, terminal escape sequences or coordinates.

\subsection{A tempting shortcut}

It is tempting to store preformatted rows as the application's state. That works until filtering needs the job identity, cancellation needs its run, or a pipe consumer needs a numeric percentage. Parsing those facts back out of painted text reverses the dependency. Here, presentation consumes the facts. Commands and tests use those same facts directly.

\subsection{Run the first build stage}

From the companion repository, run \texttt{npm ci}, then \texttt{npm run stage -- 01}. The output contains three queued jobs. Read the rows with color disabled and answer: what stable identity does each job have, what state is it in, and how much work is complete? Stage 01 introduces its own \path{stages/01-model/job.ts}, monitor snapshot and plain presentation. Run \texttt{npm run stage:check -- 01} before reading the 01-to-02 source diff. The next workshop reuses that snapshot rather than inventing a separate model for each environment.

\section{Try a real extension}

Add a job whose title contains a wide character and a combining accent. Give it an identity that remains stable if its title changes. Success means the same identity reaches every presentation and every transition. A title is a label for people; changing its spelling must not redirect an asynchronous result to another job.
